\setcounter{chapter}{2}
\chapter{Lebesgue外測度}

\paragraph{本章の目標}
Jordan外測度の「有限被覆」を「可算被覆」に置き換えて
Lebesgue外測度 $\lambda^*$ を定義し，
それが $\RR^d$ 上の外測度になることを示す．
さらに，Lebesgue外測度が開被覆によっても同じように記述できることを証明する．
%本章ではまだ可測集合を選別しない．
$\lambda^*$ はすべての部分集合に定義されるが，
一般には可算劣加法性しか言えない．

\section{動機づけ}

%第2章では，
Jordan測度は有限個の区間による近似に基づく理論であるため，
\[
A = \QQ \cap [0,1]
\]
のような集合はJordan可測にならない．
有限個の区間で $A$ を覆おうとすると，
$A$ が $[0,1]$ に稠密であるため，
結局 $[0,1]$ 全体に近い長さが必要になるからである．

しかし，もし可算個の区間を使ってよいなら事情は変わる．
$A=\{q_1,q_2,\dots\}$ と番号づけて，
各点 $q_n$ を長さ $\eps 2^{-n}$ の区間で覆えば，
全体の長さの和は
\[
\sum_{n=1}^{\infty} \eps 2^{-n} = \eps
\]
になる．
したがって，$\QQ \cap [0,1]$ の長さは $0$ とみなすのが自然である．

この観察は，Jordan理論の限界が
「被覆が有限個に制限されていること」にあることを示している．
そこで，有限被覆を可算被覆に置き換えてLebesgue外測度を導入する．

\section{拡大実数}

Lebesgue外測度では，下限や無限和を頻繁に扱う．
このとき，値域として拡大実数を設定すると記述が簡潔になって便利である．

\begin{definition}[拡大実数]
\[
\eRR := \RR \cup \{-\infty,\infty\}
\]
を拡大実数全体と呼ぶ．
特に測度論では
\[
[0,\infty]:=[0,\infty)\cup\{\infty\}
\]
を主に用いる．
\end{definition}

% \begin{remark}
% 非負の拡大実数列 $\{a_n\}_{n=1}^{\infty}$ に対して
% \[
% \sum_{n=1}^{\infty} a_n
% \]
% は部分和
% \[
% s_N:=\sum_{n=1}^N a_n
% \]
% の上限
% \[
% \sup_{N \in \NN} s_N
% \]
% として定める．
% 非負項しか扱わないので，この和は常に意味をもつ．
% （部分和列は単調増加なので，$\lim_{N \to \infty} s_N = \sup_{N \in \NN} s_N$ である．）
% \end{remark}

\begin{remark}
拡大実数では
\[
\infty-\infty,\qquad 0\cdot\infty,\qquad \frac{\infty}{\infty}
\]
のような式は一般には未定義とする．
測度論では，このような不定形を避けるため，
非負値の和・上限・下限だけを使って議論を進めることが多い．
\end{remark}

\section{Lebesgue外測度の定義}

以後，次元 $d \in \NN$ を固定する．
第2章の記法に拡大実数を含めるため，以下のように約束する
\begin{itemize}
    \item $\eRR$ の区間とは，$-\infty \le a < b \le \infty$ に対し，$I = (a,b), (a,b], [a,b), [a,b], \text{or } \emptyset$ である．便宜的に，空集合 $\emptyset$ も区間に含める．
    \item $\eRR^d$ の区間 $R = \prod_{i=1}^d I_i$ の体積を $|R| = \prod_{i=1}^d (b_i - a_i)$ で表す．また空集合の体積を $|\emptyset|:=0$ と約束する．区間に一つでも$\pm \infty$が含まれるときは $|R|=\infty$とする．
    \item $\eRR^d$ の左半開区間の全体に空集合を加えたものを
    $\calE_d := \left\{ \prod_{i=1}^d (a_i,b_i] \;\middle|\; a_i < b_i \right\} \cup \{ \emptyset\}$
%    $\calE_d := \{ \prod_{i=1}^d (a_i,b_i] \mid -\infty \le a_i < b_i \le \infty \} \cup \{ \emptyset\}$
    と書く
    \item $\eRR^d$ の基本集合（互いに素な左半開区間の有限合併）の全体を $\calA_d := \bigcup_{n \in \NN} \left\{ \bigsqcup_{j=1}^n Q_j \;\middle|\; Q_j \in \calE_d\right\}$ と書く．区間の定義に拡大実数を含めたことにより，$\calA_d$は有限加法族である．
\end{itemize}

\begin{definition}[Lebesgue外測度]
集合 $A \subset \RR^d$ に対し，$A$ のLebesgue外測度 $\lambda^*(A)$ を
\[
\lambda^*(A)
:=
\inf\left\{
\sum_{i=1}^{\infty} |R_i|
\;\middle|\;
A \subset \bigcup_{i=1}^{\infty} R_i,\ 
R_i \in \calE_d
\right\}
\]
で定める．
\end{definition}

\begin{remark}
    可算被覆の基本単位 $R_i$ として区間（$\calE_d$）ではなく基本集合（$\calA_d$）を採用しても同じ量を定める．
\end{remark}

\begin{remark}
有限被覆 $A \subset \bigcup_{i=1}^n Q_i$ は可算被覆の特別な場合
\[
A \subset \bigcup_{i=1}^\infty R_i,
\qquad
R_i :=
\begin{cases}
Q_i, & i \le n,\\
\emptyset, & i > n
\end{cases}
\]
である．
したがって，Lebesgue外測度はJordan外測度よりも精密に近似できる．
\end{remark}

\begin{proposition}[Jordan外測度との比較] \label{prop:louter-le-jouter}
有界集合 $A \subset \RR^d$ に対して
\[
\lambda^*(A) \le m_J^*(A)
\]
が成り立つ．
\end{proposition}
\begin{proof}
$A \subset E$ を満たす基本集合 $E$ を取る．
$E$ は互いに素な左半開区間の有限合併
\[
E = \bigsqcup_{j=1}^N Q_j
\]
と書け，
\[
m_J(E) = \sum_{j=1}^N |Q_j|
\]
であった．
この有限個の区間の族は $A$ の可算被覆でもあるから，
\[
\lambda^*(A) \le \sum_{j=1}^N |Q_j| = m_J(E)
\]
を得る．
$E$ の取り方について下限を取ればよい．
\end{proof}

\section{具体例}

\begin{proposition}[区間による上からの評価]
$A \subset \RR^d$ と $R \in \calE_d$ が
$A \subset R$ を満たすとする．このとき
\[
\lambda^*(A) \le |R|
\]
が成り立つ．
特に，有界集合のLebesgue外測度は有限である．
\end{proposition}
\begin{proof}
$R$ 自身による1個の被覆を考えればよい．
\end{proof}

\begin{proposition}[1点集合は零集合]
任意の $a \in \RR^d$ に対して
\[
\lambda^*(\{a\})=0
\]
が成り立つ．
\end{proposition}
\begin{proof}
非負性より $\lambda^*(\{a\}) \ge 0$ である．
任意の $\eps > 0$ に対し，$\delta > 0$ を十分小さく取って
\[
(2\delta)^d < \eps
\]
とする．
すると
\[
\{a\}
\subset
\prod_{j=1}^d (a_j-\delta,a_j+\delta)
\]
であり，右辺の体積は $(2\delta)^d$ だから
\[
\lambda^*(\{a\}) \le (2\delta)^d < \eps.
\]
$\eps$ は任意なので $\lambda^*(\{a\})=0$ である．
\end{proof}

\begin{theorem}[可算集合は零集合]
可算集合 $A \subset \RR^d$ に対して
\[
\lambda^*(A)=0
\]
が成り立つ．
\end{theorem}
\begin{proof}
$A=\{a_1,a_2,\dots\}$ と書く．
（後述の）可算劣加法性より
\[
\lambda^*(A)
\le
\sum_{n=1}^{\infty} \lambda^*(\{a_n\})
=
\sum_{n=1}^{\infty} 0
=0.
\]
非負性と合わせて $\lambda^*(A)=0$ を得る．
\end{proof}

\begin{example}
\begin{enumerate}
    \item $\QQ^d$ は可算だから $\lambda^*(\QQ^d)=0$
    \item 特に $\lambda^*(\QQ \cap [0,1]) = 0$ である
    \item $\ZZ^d$ も可算だから $\lambda^*(\ZZ^d)=0$
\end{enumerate}
\end{example}

% \begin{remark}
% Jordan理論の範囲では，
% Dirichlet関数の台集合$A=\QQ \cap [0,1]$や
% Lebesgue外測度では長さ $0$ になる．
% 非有界集合もJordan可測でなかったが，
% Lebesgue
% 有限被覆と可算被覆の顕著な違いである．
% \end{remark}

\section{Lebesgue外測度の性質}

\begin{definition}[Carathéodory外測度]
集合 $X$ 上の集合関数
\[
\mu^* : \mathcal P(X) \to [0,\infty]
\]
が（Carathéodory）外測度であるとは，次の3条件を満たすことをいう．
\begin{enumerate}
    \item $\mu^*(\emptyset)=0$
    \item （単調性）$A \subset B$ ならば $\mu^*(A) \le \mu^*(B)$
    \item （可算劣加法性）任意の集合列 $\{A_n\}_{n=1}^{\infty}$ に対して
    \[
    \mu^*\left(\bigcup_{n=1}^{\infty} A_n\right)
    \le
    \sum_{n=1}^{\infty} \mu^*(A_n)
    \]
    が成り立つ
\end{enumerate}
\end{definition}

\begin{theorem}[Lebesgue外測度は外測度]
$\lambda^*$ は $\RR^d$ 上のCarathéodory外測度である．
\end{theorem}

% \begin{theorem}[Lebesgue外測度の基本性質]
% 任意の $A,B,A_n \subset \RR^d$ に対して次が成り立つ．
% \begin{enumerate}
%     \item  $\lambda^*(\emptyset)=0$
%     \item （単調性） $A \subset B$ ならば $\lambda^*(A)\le \lambda^*(B)$
%     \item （可算劣加法性）
%     \[
%     \lambda^*\left(\bigcup_{n=1}^{\infty} A_n\right)
%     \le
%     \sum_{n=1}^{\infty}\lambda^*(A_n)
%     \]
% \end{enumerate}
% \end{theorem}

\begin{proof}
\begin{enumerate}
    \item $\lambda^*(\emptyset)=0$：
    まず定義から $\lambda^*(\emptyset)\ge 0$ である．
    一方，
    \[
    \emptyset \subset \bigcup_{n=1}^{\infty}\emptyset
    \]
    であるから
    \[
    \lambda^*(\emptyset) \le \sum_{n=1}^{\infty} |\emptyset| = 0.
    \]
    よって $\lambda^*(\emptyset)=0$ である．

    \item 単調性：
    $A \subset B$ とする．
    $B$ の任意の可算被覆
    \[
    B \subset \bigcup_{n=1}^{\infty} R_n
    \]
    はそのまま $A$ の可算被覆でもある．
    したがって，$A$ に対する下限は $B$ に対する下限以下であり，
    \[
    \lambda^*(A) \le \lambda^*(B)
    \]
    を得る．

    \item 可算劣加法性：
    集合列 $\{A_n\}_{n=1}^{\infty}$ を取る．
    もし
    \[
    \sum_{n=1}^{\infty}\lambda^*(A_n)=\infty
    \]
    なら不等式は自明である．
    したがって右辺が有限であるとし，$\eps > 0$ を任意に取る．
    各 $n$ に対して，$\lambda^*(A_n)$ の定義より
    $A_n$ を覆う可算個の区間 $\{R_{n,k}\}_{k=1}^{\infty}$ を
    \[
    A_n \subset \bigcup_{k=1}^{\infty} R_{n,k},
    \qquad
    \sum_{k=1}^{\infty} |R_{n,k}|
    <
    \lambda^*(A_n) + \frac{\eps}{2^n}
    \]
    となるように取れる．
    すると二重列 $\{R_{n,k}\}_{n,k}$ は可算個の区間からなり，
    \[
    \bigcup_{n=1}^{\infty} A_n
    \subset
    \bigcup_{n=1}^{\infty}\bigcup_{k=1}^{\infty} R_{n,k}
    \]
    であるから，
    \begin{align*}
    \lambda^*\left(\bigcup_{n=1}^{\infty} A_n\right)
    &\le
    \sum_{n=1}^{\infty}\sum_{k=1}^{\infty} |R_{n,k}| \\
    &<
    \sum_{n=1}^{\infty}
    \left(
    \lambda^*(A_n)+\frac{\eps}{2^n}
    \right) \\
    &=
    \sum_{n=1}^{\infty}\lambda^*(A_n)+\eps.
    \end{align*}
    $\eps > 0$ は任意だから
    \[
    \lambda^*\left(\bigcup_{n=1}^{\infty} A_n\right)
    \le
    \sum_{n=1}^{\infty}\lambda^*(A_n)
    \]
    が従う．
\end{enumerate}
\end{proof}

% \begin{corollary}[有限劣加法性]
% 任意の $A,B \subset \RR^d$ に対して
% \[
% \lambda^*(A \cup B) \le \lambda^*(A) + \lambda^*(B)
% \]
% が成り立つ．
% \end{corollary}
% \begin{proof}
% 上の定理で
% \[
% A_1=A,\qquad A_2=B,\qquad A_n=\emptyset\ (n \ge 3)
% \]
% とおけばよい．
% \end{proof}

\begin{proposition}[平行移動不変性]
任意の $A \subset \RR^d$ と $c \in \RR^d$ に対して
\[
\lambda^*(A+c)=\lambda^*(A)
\]
が成り立つ．
\end{proposition}
\begin{proof}
$A \subset \bigcup_{n=1}^{\infty} R_n$ ならば
\[
A+c \subset \bigcup_{n=1}^{\infty} (R_n+c)
\]
であり，平行移動は各辺の長さを変えないので
\[
|R_n+c| = |R_n|
\]
である．
したがって
\[
\lambda^*(A+c) \le \sum_{n=1}^{\infty} |R_n|
\]
となるから，$A$ の被覆について下限を取って
\[
\lambda^*(A+c) \le \lambda^*(A)
\]
を得る．
逆向きの不等式は $-c$ だけ平行移動すれば従う．
\end{proof}

\section{Jordan測度との整合}

\begin{theorem}[{$m_J : \calA_J \to [0,\infty]$の拡張}]
    Jordan可測集合 $A$ に対し，
    \[
    m_J(A) = \lambda^*(A).
    \]
\end{theorem}

\begin{proof}
Jordan可測集合 $A$ は有界であり，$m_J(A) = m_J^*(A) = m_{J,*}(A) <\infty$ である．
$A = \emptyset$ のとき $m_J(\emptyset)=0, \lambda^*(\emptyset)=0$ より成り立つ．
以下では$A \neq \emptyset$ とする．
まず $\lambda^*(A)\le m_J(A)$ は，\cref{prop:louter-le-jouter} から直ちに従う．
% 任意の $\eps>0$ に対し，$m_J(A)=m_J^*(A)$ とJordan外測度の定義より，
% 基本集合 $G\in\calA_d$ が存在して
% \[
% A\subset G,
% \qquad
% m_J(G)<m_J(A)+\eps
% \]
% を満たす．
% $G=\bigsqcup_{j=1}^N Q_j$ と半開区間の互いに素な有限合併に書けば，
% $\{Q_j\}_{j=1}^N$ は $A$ の可算被覆とみなせるので
% \[
% \lambda^*(A)
% \le
% \sum_{j=1}^N |Q_j|
% =
% m_J(G)
% <
% m_J(A)+\eps .
% \]
% $\eps>0$ は任意だから $\lambda^*(A)\le m_J(A)$ である．

逆向きの不等式を示す．
$A\subset\bigcup_{n=1}^\infty R_n$ を満たす任意の可算被覆
$\{R_n\}_{n=1}^\infty\subset\calE_d$ を取る．
$\sum_{n=1}^\infty |R_n|=\infty$ の場合は自明なので，
以下ではこの和が有限であるとする．
このとき各 $R_n$ は有界である．

基本集合 $F\subset A$ を任意に取る．
任意の $\eps>0$ に対し，各 $R_n$ を少し膨らませた開区間 $U_n$ を
\[
R_n\subset U_n,
\qquad
|U_n|\le |R_n|+\frac{\eps}{2^n}
\]
となるように取る．
%
% $F=\bigsqcup_{j=1}^N Q_j$，
% $Q_j=\prod_{i=1}^d(a_{ij},b_{ij}]$ と書くと，$F\subset A$ だから各端点は有限である．
% 十分小さい $\eta>0$ に対して
% \[
% K_\eta
% :=
% \bigsqcup_{j=1}^N \prod_{i=1}^d [a_{ij}+\eta,b_{ij}]
% \subset F
% \]
% はコンパクトなJordan可測集合であり，直接計算により
% $m_J(K_\eta)\to m_J(F)$ $(\eta\downarrow0)$ である．
% したがって任意の $\delta>0$ に対し，$\eta>0$ を十分小さく取れば
% \[
% m_J(F)-\delta\le m_J(K_\eta)
% \]
% となる．
さらに任意の $\delta >0$ に対し，$F$を少し縮ませたコンパクト集合 $K_\delta$ を
\[
K_\delta \subset F, \qquad m_J(F) - \delta \le m_J(K_\delta)
\]
となるように取る．
%
$\{U_n\}_{n=1}^\infty$ はコンパクト集合 $K_\delta$ の開被覆だから有限部分被覆を持つ．
Jordan外測度の有限劣加法性より
\[
m_J(K_\delta)=m_J^*(K_\delta)
\le
\sum_{k=1}^M m_J^*(U_{n_k})
\le
\sum_{k=1}^M |U_{n_k}|
\le
\sum_{n=1}^\infty |R_n|+\eps .
\]
ゆえに
\[
m_J(F)-\delta
\le
\sum_{n=1}^\infty |R_n|+\eps .
\]
$\delta,\eps>0$ は任意なので
\[
m_J(F)\le \sum_{n=1}^\infty |R_n|
\]
である．
最後に $F\subset A$ を満たす基本集合全体について上限を取り，
% $A$ のJordan可測性より
% \[
% m_J(A)=m_{J,*}(A)
% \le
% \sum_{n=1}^\infty |R_n|
% \]
% を得る．
可算被覆について下限を取れば \[m_J(A)=m_{J,*}(A)\le\lambda^*(A)\] を得る．
以上より $m_J(A)=\lambda^*(A)$ である．
\end{proof}

\begin{theorem}[$\calA_J$上で可算加法的]
    互いに素なJordan可測集合の列 $\{A_i\}_{i=1}^\infty$ に対し，
    \[
    \lambda^*\left( \bigsqcup_{i=1}^\infty A_i \right) = \sum_{i=1}^\infty \lambda^*(A_i).
    \]
\end{theorem}

\begin{proof}
$B_n:=\bigsqcup_{i=1}^n A_i$ とおく．
各 $B_n$ は有限個のJordan可測集合の互いに素な合併だからJordan可測であり，
Jordan測度の有限加法性と前定理より
\[
\lambda^*(B_n)=m_J(B_n)=\sum_{i=1}^n m_J(A_i)=\sum_{i=1}^n \lambda^*(A_i)
\]
が成り立つ．

いま
\[
B_n \subset \bigsqcup_{i=1}^\infty A_i
\]
だから単調性より
\[
\sum_{i=1}^n \lambda^*(A_i)=\lambda^*(B_n)
\le
\lambda^*\left( \bigsqcup_{i=1}^\infty A_i \right)
\]
である．$n\to\infty$ とすれば
\[
\sum_{i=1}^\infty \lambda^*(A_i)
\le
\lambda^*\left( \bigsqcup_{i=1}^\infty A_i \right)
\]
を得る．

逆向きの不等式はLebesgue外測度の可算劣加法性から
\[
\lambda^*\left( \bigsqcup_{i=1}^\infty A_i \right)
\le
\sum_{i=1}^\infty \lambda^*(A_i)
\]
である．
以上より等号が成り立つ．
\end{proof}

\section{Lebesgue外測度の外正則性}%{開被覆による書き換え}

$\RR^d$の開集合系を$\calO_d := \{ O \subset \RR^d \mid O \text{ は開集合}\}$と書く．
\begin{remark}
    一般の集合$X$の部分集合族$\calO_X$が開集合系であるとは，以下の3条件を満たすことをいう
    \begin{itemize}
        \item $X,\emptyset \in \calO_X$
        \item （任意合併）$\{ O_\lambda \}_{\lambda \in \Lambda} \subset \calO \implies \bigcup_{\lambda \in \Lambda} O_\lambda \in \calO$
        \item （有限交叉）$O_1,O_2 \in \calO_X \implies O_1 \cap O_2 \in \calO_X$
    \end{itemize}
    特に，$\RR^d$の開集合系$\calO_d$は，上記の条件を満たす．
\end{remark}

% \begin{theorem}[開区間被覆による書き換え]
% 任意の $A \subset \RR^d$ に対して
% \[
% \lambda^*(A)
% =
% \inf\left\{
% \sum_{n=1}^{\infty} |U_n|
% \;\middle|\;
% A \subset \bigcup_{n=1}^{\infty} U_n,
% U_n \text{ は有界開区間}
% \right\}
% \]
% が成り立つ．
% \end{theorem}
% \begin{proof}
% 右辺を $\alpha(A)$ とおく．
% \[
% A \subset \bigcup_{n=1}^{\infty} U_n,
% \qquad
% U_n \text{ は有界開区間}
% \]
% とする．
% 各 $n$ に対し，任意の $\delta>0$ に対して
% $U_n$ を含む左半開区間 $R_n \in \calE_d$ を
% \[
% |R_n| < |U_n|+\delta
% \]
% となるように取れるから，
% \[
% \lambda^*(U_n)\le |U_n|
% \]
% である．
% したがって可算劣加法性より
% \[
% \lambda^*(A)
% \le
% \sum_{n=1}^{\infty}\lambda^*(U_n)
% \le
% \sum_{n=1}^{\infty}|U_n|
% \]
% を得る．
% 被覆について下限を取れば $\lambda^*(A)\le \alpha(A)$ である．
%
% 逆向きの不等式を示す．
% もし $\lambda^*(A)=\infty$ なら，
% すでに示した $\lambda^*(A)\le \alpha(A)$ より
% $\alpha(A)=\infty$ であり，自明である．
%
% $\eps>0$ を取る．
% Lebesgue外測度の定義より，
% $A$ を覆う区間列 $\{R_n\}_{n=1}^{\infty} \subset \calE_d$ で
% \[
% \sum_{n=1}^{\infty}|R_n|<\lambda^*(A)+\eps
% \]
% を満たすものが取れる．
% 右辺は有限だから各 $|R_n|$ は有限であり，したがって各 $R_n$ は有界である．
% 各 $n$ に対して，$R_n \subset U_n$ を満たす有界開区間 $U_n$ を
% \[
% |U_n|<|R_n|+\frac{\eps}{2^n}
% \]
% となるように選ぶ．
% すると $\{U_n\}$ は $A$ の有界開区間被覆であり，
% \begin{align*}
% \sum_{n=1}^{\infty}|U_n|
% &<
% \sum_{n=1}^{\infty}\left(|R_n|+\frac{\eps}{2^n}\right) \\
% &=
% \sum_{n=1}^{\infty}|R_n|+\eps \\
% &<
% \lambda^*(A)+2\eps.
% \end{align*}
% したがって
% \[
% \alpha(A)\le \lambda^*(A)+2\eps
% \]
% を得る．
% $\eps>0$ は任意だから $\alpha(A)\le \lambda^*(A)$ である．
% \end{proof}
%
% \begin{corollary}[開集合による外正則性]
% 任意の $A \subset \RR^d$ に対して
% \[
% \lambda^*(A)
% =
% \inf\{\lambda^*(G)\mid A \subset G,\ G \text{ は開集合}\}
% \]
% が成り立つ．
% \end{corollary}
% \begin{proof}
% 右辺を $\beta(A)$ とおく．
% $A \subset G$ なら単調性より
% \[
% \lambda^*(A)\le \lambda^*(G)
% \]
% だから
% \[
% \lambda^*(A)\le \beta(A)
% \]
% である．
%
% 逆向きの不等式を示す．
% $\eps>0$ を取る．
% 前定理より，$A$ を覆う有界開区間列 $\{U_n\}_{n=1}^{\infty}$ で
% \[
% \sum_{n=1}^{\infty}|U_n|<\lambda^*(A)+\eps
% \]
% を満たすものが取れる．
% \[
% G:=\bigcup_{n=1}^{\infty}U_n
% \]
% とおけば，$G$ は開集合で $A \subset G$ である．
% また外測度の可算劣加法性より
% \[
% \lambda^*(G)
% \le
% \sum_{n=1}^{\infty}\lambda^*(U_n)
% \le
% \sum_{n=1}^{\infty}|U_n|
% <
% \lambda^*(A)+\eps.
% \]
% ここで第2の不等式は，前定理を $A=U_n$ に適用すれば
% $\lambda^*(U_n)\le |U_n|$ と分かることによる．
% したがって
% \[
% \beta(A)\le \lambda^*(A)+\eps.
% \]
% $\eps>0$ は任意だから $\beta(A)\le \lambda^*(A)$ である．
% \end{proof}

\begin{theorem}[開集合による外正則性]
任意の $A \subset \RR^d$ に対して
\[
\lambda^*(A)
=
\inf\{\lambda^*(G)\mid A \subset G,\ G \text{ は開集合}\}
\]
が成り立つ．
\end{theorem}

\begin{proof}
    右辺を $\omega(A) := \inf\{ \lambda^*(O) \mid A \subset O, O \in \calO_d \}$ とおく．
    まずLebesgue外測度の単調性により，任意の開集合 $O$ に対し \[A \subset O \implies \lambda^*(A) \le \lambda^*(O)\] である．よって $\lambda^*(A) \le \omega(A)$ である．

    逆向きの不等号を示す．
    もし $\lambda^*(A)=\infty$ なら，すでに示した $\lambda^*(A)\le\omega(A)$ より
    $\omega(A)=\infty$ であり，自明である．

    以下，$\lambda^*(A)<\infty$ とする．
    $\eps>0$ を取る．
    Lebesgue外測度の定義より，
    区間列 $\{ Q_i \}_{i=1}^\infty \subset \calE_d$ が存在して
    \[
    A \subset \bigcup_{i=1}^\infty Q_i,
    \qquad
    \sum_{i=1}^\infty |Q_i| < \lambda^*(A) + \frac{\eps}{2}
    \]
    である．このとき右辺は有限なので，空でない $Q_i$ は有界である．
    各 $i$ に対し，$Q_i\subset O_i$ を満たす開集合 $O_i\in\calO_d$ を
    \[
    \lambda^*(O_i) < |Q_i| + \frac{\eps}{2^{i+1}}
    \]
    となるように取る．
    実際，$Q_i=\emptyset$ なら $O_i=\emptyset$ とし，
    $Q_i\neq\emptyset$ なら，まず $Q_i\subset V_i$ かつ
    \[
    |V_i|<|Q_i|+\frac{\eps}{2^{i+2}}
    \]
    となる有界開区間 $V_i$ を取る．
    次に $V_i\subset R_i$ かつ
    \[
    |R_i|<|V_i|+\frac{\eps}{2^{i+2}}
    \]
    となる左半開区間 $R_i\in\calE_d$ を取り，$O_i:=V_i$ とおけば
    \[
    \lambda^*(O_i)\le |R_i|
    <
    |Q_i|+\frac{\eps}{2^{i+1}}
    \]
    である．

    したがって，$O_\eps := \bigcup_{i=1}^\infty O_i$ とおくと，
    $O_\eps$ は $A$ を含む開集合（$A \subset O_\eps, O_\eps \in \calO_d$）であって，
    \[
    \lambda^*(O_\eps)
    \le
    \sum_{i=1}^\infty \lambda^*(O_i)
    <
    \sum_{i=1}^\infty |Q_i| + \frac{\eps}{2}
    <
    \lambda^*(A) + \eps
    \]
    である．$\eps>0$ は任意だから $\omega(A) \le \lambda^*(A)$ である．
\end{proof}

\section{解釈}

%現代的には，
Lebesgue理論の出発点は
「すべての部分集合に可算加法的な体積を与えること」ではなく，
まず外測度
\[
\lambda^* : \mathcal P(\RR^d) \to [0,\infty]
\]
を全ての部分集合上に定めることである．
本章で得た $\lambda^*$ は
\begin{enumerate}
    \item 全域的に定義される
    \item 非負で，$\lambda^*(\emptyset)=0$ を満たす
    \item 可算劣加法的である
    \item 平行移動不変である
    \item 区間の体積や可算集合の長さ $0$ という幾何学的直観と整合する
\end{enumerate}
という性質をもつ．
ただし，この段階ではまだ可算加法性は一般には成り立たないので，
$\lambda^*$ は「測度」ではなく「外測度」にとどまっている．

%したがって $\lambda^*$ は，まだ「測度」ではなく「外測度」である．
次章以降ではこれに対応する内測度を導入し，
\[
\text{外測度}=\text{内測度}
\]
が成り立つ集合をLebesgue可測集合として定義する．
